\documentclass[../../main.tex]{subfiles}
\begin{document}
% SOURCE p=254; printed=239; object=ch08-preliminaries-exercises-title; region=upper-middle
\Needspace{12\baselineskip}
\subsection*{习题}
% SOURCE p=254; printed=239; object=ch08-ex01-leibniz; number=1; region=middle
\begin{exercise}\label{exercise:ch08-leibniz}
证明正文中给出的偏导数乘积法则。由此推出
\begin{align*}
\partial^\alpha(x^\beta f)&=x^\beta\partial^\alpha f+
\sum_{\gamma,\delta}c_{\gamma\delta}x^\delta\partial^\gamma f,\\
x^\beta\partial^\alpha f&=\partial^\alpha(x^\beta f)+
\sum_{\gamma,\delta}c'_{\gamma\delta}\partial^\gamma(x^\delta f),
\end{align*}
其中系数为常数，而且除非 $|\gamma|<|\alpha|$ 且 $|\delta|<|\beta|$，否则 $c_{\gamma\delta}=c'_{\gamma\delta}=0$。
\end{exercise}
% SUPPLEMENT: id=v14-ch08-ex01-solution; parent=ch08-ex01-leibniz; kind=solution+motivation+hypotheses; origin=independent-derivation
\begin{solution}
\textbf{完整补写（重构新增）。}先对一个坐标证明
\begin{align*}
\partial_j^a(fg)=\sum_{b=0}^a\binom ab(\partial_j^bf)(\partial_j^{a-b}g).
\end{align*}
$a=0$ 显然。若对 $a$ 成立，再求一次 $\partial_j$，每项分成导数落在第一或第二因子的两项；相同项的系数相加为
\begin{align*}
\binom a{b-1}+\binom ab=\binom{a+1}b.
\end{align*}
边界处不存在的系数取零，即得 $a+1$ 的公式。依次在 $n$ 个坐标应用该式，连续混合偏导可交换，系数相乘为
$\prod_j\binom{\alpha_j}{\beta_j}=\alpha!/(\beta!(\alpha-\beta)!)$，从而得到\cref{proposition:ch08-leibniz-formula}。

对单项式求导，若 $\rho\leqslant\beta$，则 $\partial^\rho x^\beta=\beta!x^{\beta-\rho}/(\beta-\rho)!$；否则为零。因此乘积法则给出具体的有限展开
\begin{align*}
\partial^\alpha(x^\beta f)=
\sum_{\substack{\rho\leqslant\alpha\\\rho\leqslant\beta}}
\binom\alpha\rho\frac{\beta!}{(\beta-\rho)!}
 x^{\beta-\rho}\partial^{\alpha-\rho}f,
\qquad \binom\alpha\rho=\frac{\alpha!}{\rho!(\alpha-\rho)!}.
\end{align*}
$\rho=0$ 给主项；其余 $\rho\ne0$，令 $\gamma=\alpha-\rho$、$\delta=\beta-\rho$，则 $|\gamma|<|\alpha|$、$|\delta|<|\beta|$，这就证明第一式及系数条件。

第二式按 $|\alpha|$ 归纳，并在每阶同时允许任意 $\beta$。$|\alpha|=0$ 时没有余项。若较低阶已证，把第一式解成
\begin{align*}
x^\beta\partial^\alpha f=\partial^\alpha(x^\beta f)
-\sum_{\gamma,\delta}c_{\gamma\delta}x^\delta\partial^\gamma f.
\end{align*}
对右侧每个低阶项应用归纳假设，将其写成 $\partial^\gamma(x^\delta f)$ 与有限多个 $\partial^{\gamma'}(x^{\delta'}f)$ 的和；后者满足 $|\gamma'|<|\gamma|$、$|\delta'|<|\delta|$。所有项因此仍同时满足低于原来 $|\alpha|$、$|\beta|$ 的严格不等式。有限次相加并收集同类项，就得所求常系数 $c'_{\gamma\delta}$。若 $\beta=0$，第一式无低阶项，第二式也无低阶项，边界情形包括在内。
\end{solution}
\begin{remark}
\textbf{动机与条件（重构新增）。}求导与乘单项式不交换，但交换产生的误差在微分阶数和多项式次数两方面都下降；这使三角形归纳可行，也是 Schwartz 空间两种刻画等价的关键。固定 $\alpha$ 只需相应阶连续偏导数；常数系数仅由 $\alpha,\beta$ 决定，与 $x,f$ 无关。两个严格降阶条件都已由 $\rho\ne0$ 或归纳传递验证，不能只证明其中一个。
\end{remark}
% END SUPPLEMENT: id=v14-ch08-ex01-solution
% SOURCE p=254; printed=239; object=ch08-ex02-multinomial; number=2; region=middle
\begin{exercise}\label{exercise:ch08-multinomial}
注意二项式定理可写成
\begin{align*}
(x_1+x_2)^k=\sum_{|\alpha|=k}\frac{k!}{\alpha!}x^\alpha
\qquad\bigl(x=(x_1,x_2),\ \alpha=(\alpha_1,\alpha_2)\bigr).
\end{align*}
证明以下推广。
\begin{enumerate}
\item\label{item:ch08-multinomial-a}多项式定理：若 $x\in\mathbb R^n$，则
$\displaystyle (x_1+\cdots+x_n)^k=\sum_{|\alpha|=k}\frac{k!}{\alpha!}x^\alpha$。
\item\label{item:ch08-multinomial-b}$n$ 维二项式定理：若 $x,y\in\mathbb R^n$，则
$\displaystyle (x+y)^\alpha=\sum_{\beta+\gamma=\alpha}\frac{\alpha!}{\beta!\gamma!}x^\beta y^\gamma$。
\end{enumerate}
\end{exercise}
% SUPPLEMENT: id=v14-ch08-ex02-solution; parent=ch08-ex02-multinomial; kind=solution+motivation+hypotheses; origin=independent-derivation
\begin{solution}
\textbf{完整补写（重构新增）。}(a) 对非负整数 $k$，将 $k$ 个因子 $(x_1+\cdots+x_n)$ 的乘积按分配律展开。每项由在每个因子中选一个坐标确定。若共选 $x_j$ 恰 $\alpha_j$ 次，则 $|\alpha|=k$，所得单项式为 $x^\alpha$。这样的选择共有
\begin{align*}
\binom{k}{\alpha_1}\binom{k-\alpha_1}{\alpha_2}\cdots
\binom{k-\alpha_1-\cdots-\alpha_{n-1}}{\alpha_n}
=\frac{k!}{\alpha_1!\cdots\alpha_n!}
\end{align*}
种：依次选择各坐标出现的位置。收集相同单项式即得结论。$k=0$ 时左侧为空乘积 $1$，右侧只有 $\alpha=0$ 一项，也是 $1$。

(b) 对每个坐标使用一元二项式定理，再相乘：
\begin{align*}
(x+y)^\alpha
&=\prod_{j=1}^n(x_j+y_j)^{\alpha_j}
=\prod_{j=1}^n\sum_{\beta_j=0}^{\alpha_j}
\frac{\alpha_j!}{\beta_j!(\alpha_j-\beta_j)!}
 x_j^{\beta_j}y_j^{\alpha_j-\beta_j}\\
&=\sum_{\beta+\gamma=\alpha}
\frac{\alpha!}{\beta!\gamma!}x^\beta y^\gamma.
\end{align*}
所有和与乘积都有限，分配律即可交换其顺序。某个 $\alpha_j=0$ 时该因子只有一项，仍成立。
\end{solution}
\begin{remark}
\textbf{动机与条件（重构新增）。}(a) 按总次数分配多个坐标，(b) 在每个坐标内分别分配给两个向量；两种计数解释了 Taylor 公式和 Leibniz 公式中不同的阶乘系数。非负整数阶确保有限展开；实数条件足够，事实上复数或任意交换代数中的元素也可作同样的有限运算，但本题不需要该推广。
\end{remark}
% END SUPPLEMENT: id=v14-ch08-ex02-solution
% SOURCE p=254; printed=239; object=ch08-ex03-flat-function; number=3; region=lower
\begin{exercise}\label{exercise:ch08-flat-function}
令 $\eta(t)=e^{-1/t}$（$t>0$），$\eta(t)=0$（$t\leqslant0$）。
\begin{enumerate}
\item\label{item:ch08-flat-a}若 $k\in\mathbb N$、$t>0$，则 $\eta^{(k)}(t)=P_k(1/t)e^{-1/t}$，其中 $P_k$ 为 $2k$ 次多项式。
\item\label{item:ch08-flat-b}对每个 $k\in\mathbb N$，$\eta^{(k)}(0)$ 均存在且等于零。
\end{enumerate}
\end{exercise}
% SUPPLEMENT: id=v14-ch08-ex03-solution; parent=ch08-ex03-flat-function; kind=solution+motivation+hypotheses; origin=independent-derivation
\begin{solution}
\textbf{完整补写（重构新增）。}(a) 可从 $k=0$、$P_0(s)=1$ 开始归纳。令 $s=1/t$，则 $\mathrm ds/\mathrm dt=-s^2$。若公式对 $k$ 成立，求导得
\begin{align*}
\eta^{(k+1)}(t)=s^2\bigl(P_k(s)-P_k'(s)\bigr)e^{-s}.
\end{align*}
因此取 $P_{k+1}(s)=s^2(P_k(s)-P_k'(s))$。归纳地，$P_k$ 首项系数为 $1$、次数为 $2k$；求导项次数较低，不会消去最高次项，故 $P_{k+1}$ 的次数恰为 $2k+2$。

(b) 先说明对任何非负整数 $d$，$s^de^{-s}\to0$（$s\to\infty$）：选整数 $l>d$，由 $e^s\geqslant s^l/l!$ 得 $s^de^{-s}\leqslant l!s^{d-l}\to0$。因此任意多项式乘 $e^{-s}$ 都趋零。首先 $\eta(t)\to0=\eta(0)$（$t\downarrow0$），而左侧恒为零，所以 $\eta$ 连续。假设 $\eta$ 已为 $C^k$，且零点各阶导数均为零。其第 $k$ 阶导数在正半轴为 (a) 中的表达式，在负半轴为零。零点的右差商是
\begin{align*}
\frac{\eta^{(k)}(t)-\eta^{(k)}(0)}{t}
=(1/t)P_k(1/t)e^{-1/t}\longrightarrow0,
\end{align*}
左差商为零，故 $\eta^{(k+1)}(0)=0$。正半轴上的 $\eta^{(k+1)}(t)=P_{k+1}(1/t)e^{-1/t}\to0$，负半轴仍为零，说明这个导数也在零点连续。由归纳，$\eta\in C^\infty$ 且所有零点导数均为零。
\end{solution}
\begin{remark}
\textbf{动机与条件（重构新增）。}指数的衰减压过求导带来的任意有限次 $1/t$ 增长，因而可以光滑地接上零函数。只证明正半轴各阶导数趋零还不等于已经证明零点可微；上面的差商归纳补足了这一环节。$t>0$ 是表达式 $e^{-1/t}$ 的分支条件，$t\leqslant0$ 的零延拓则是构造平坦边界的关键。函数在零点并不由其 Taylor 级数恢复：全部系数为零，但它在每个正点严格为正，这也说明光滑性不同于实解析性。
\end{remark}
% END SUPPLEMENT: id=v14-ch08-ex03-solution
% SOURCE p=254; printed=239; object=ch08-ex04-linfty-uc; number=4; region=lower
\begin{exercise}\label{exercise:ch08-linfty-uc-representative}
若 $f\in L^\infty$ 且 $\|\tau_yf-f\|_\infty\to0$（$y\to0$），则 $f$ 几乎处处等于一个一致连续函数。（令 $A_rf$ 如\cref{theorem:ch03-lebesgue-differentiation}中所定义。则 $r>0$ 时 $A_rf$ 一致连续，且 $r\to0$ 时这些函数在一致范数下为 Cauchy。）
\end{exercise}
% SUPPLEMENT: id=v14-ch08-ex04-solution; parent=ch08-ex04-linfty-uc; kind=solution+motivation+hypotheses; origin=independent-derivation
\begin{solution}
\textbf{完整补写（重构新增）。}先选取处处有界的 Borel 代表，仍记为 $f$，使 $|f(x)|\leqslant M=\|f\|_\infty$。这样的代表由 Lebesgue 测度的完备化版本定理得到，再在超出本质界的 Borel 零集上改成零即可。改动代表不改变任何平移的 $L^\infty$ 等价类。对 $r>0$ 定义
\begin{align*}
A_rf(x)=\frac1{m(B_r)}\int_{B_r(x)}f(u)\,\mathrm du
=\frac1{m(B_r)}\int_{B_r}f(x-y)\,\mathrm dy,
\qquad B_r=B(0,r).
\end{align*}
这在每个 $x$ 处有定义，且与原函数在零集上的取值无关，$|A_rf(x)|\leqslant M$。又有
\begin{align*}
|A_rf(x+h)-A_rf(x)|
\leqslant\frac{M}{m(B_r)}m\bigl((B_r+h)\mathbin\triangle B_r\bigr).
\end{align*}
若 $|h|<r$，右侧的对称差包含于球壳 $B(0,r+|h|)\setminus B(0,r-|h|)$（边界不影响测度），其测度至多
\begin{align*}
m(B_1)\bigl((r+|h|)^n-(r-|h|)^n\bigr)\longrightarrow0.
\end{align*}
上界与 $x$ 无关，故每个 $A_rf$ 一致连续。

令 $\omega(r)=\sup_{|y|<r}\|\tau_yf-f\|_\infty$。它有限且由题设趋于零。对每个固定的 $y\in B_r$，$|f(x-y)-f(x)|\leqslant\omega(r)$ 对几乎所有 $x$ 成立；但不能直接对不可数个 $y$ 的例外零集取并。为严谨处理，集合
\begin{align*}
D_r=\{(x,y)\in\mathbb R^n\times B_r:
|f(x-y)-f(x)|>\omega(r)\}
\end{align*}
为 Borel 集，每个固定 $y$ 的 $x$ 截面测度为零。对非负示性函数应用 Tonelli 定理，得乘积测度为零，故对几乎所有 $x$，不等式对几乎所有 $y\in B_r$ 成立。于是
\begin{align*}
|A_rf(x)-f(x)|\leqslant
\frac1{m(B_r)}\int_{B_r}|f(x-y)-f(x)|\,\mathrm dy
\leqslant\omega(r)
\end{align*}
对几乎所有 $x$ 成立。因此 $\|A_rf-f\|_\infty\leqslant\omega(r)$。

连续函数的一致范数等于其本质上确界：若某点的绝对值严格大于一个界，连续性使某个非空开球内仍严格大于该界，而该球有正测度。由于 $A_rf-A_sf$ 连续，得到
\begin{align*}
\|A_rf-A_sf\|_u=\|A_rf-A_sf\|_\infty
\leqslant\omega(r)+\omega(s)\longrightarrow0\qquad(r,s\downarrow0).
\end{align*}
例如取 $r_j=1/j$，这些有界连续函数在一致范数下为 Cauchy，故一致收敛到一个有界函数 $g$。$g$ 一致连续：给定 $\varepsilon>0$，选 $j$ 使 $\|g-A_{r_j}f\|_u<\varepsilon/3$，再利用 $A_{r_j}f$ 的一致连续性控制中间差，三角不等式即得所需的统一 $\delta$。最后
\begin{align*}
\|g-f\|_\infty\leqslant
\|g-A_{r_j}f\|_u+\omega(r_j)\longrightarrow0,
\end{align*}
所以 $g=f$ 几乎处处。也可用 Lebesgue 微分定理识别同一个极限；这里的估计已经直接完成了识别。
\end{solution}
\begin{remark}
\textbf{动机与条件（重构新增）。}局部球平均先把有界可测函数变成一致连续函数，平移连续性再保证平均与原函数在本质一致范数中接近。有界性控制球移动时的对称差；$r>0$ 保证球体积正且有限；Borel 代表使两变量函数可直接应用 Tonelli；Lebesgue 测度的平移不变性及非空开集正测度分别用于平均公式和两种上确界的比较。几乎处处相等是自然结论，因为任意改动一个零集上的值都不改变题设，却可能破坏逐点连续性。这个一致连续代表唯一：两个连续代表的差若某点非零，就在一个正测度开球上非零，不能几乎处处为零。反过来，若 $f$ 有有界一致连续代表 $g$，则 $\|\tau_yf-f\|_\infty\leqslant\|\tau_yg-g\|_u\to0$；所以本题刻画了 $L^\infty$ 中平移连续的元素。
\end{remark}
% END SUPPLEMENT: id=v14-ch08-ex04-solution
\end{document}
